% answering.tex -- Module 7: Answering HSC questions.
% Key verbs (NESA glossary of key words, explained in simple words) and model answers
\unnumberedchapter[\faPenNib]{Answering HSC questions}\label{bm:answering}

Every HSC question starts with a \textbf{key verb}, such as \emph{explain} or
\emph{evaluate}. The verb tells you what kind of answer to write. The marks tell you how much
to write: plan about one clear point for each mark.

\section*{What each key verb asks for}

{\small
\renewcommand\arraystretch{1.2}\arrayrulecolor{hscrule}%
\setlength\LTpre{4pt}\setlength\LTpost{4pt}%
\begin{longtable}{@{}L{18mm}@{\hspace{3mm}}L{66mm}@{\hspace{3mm}}L{60mm}@{}}
  \colhead{Verb} & \colhead{What you must do} & \colhead{A Module 7 example}\\\hline
  \endfirsthead
  \colhead{Verb} & \colhead{What you must do} & \colhead{A Module 7 example}\\\hline
  \endhead
  \textbf{Identify} & Recognise something and name it. A word or a phrase is enough. &
    Identify the type of spectrum given by a hot filament.\\\hline
  \textbf{Outline} & Give the main features only. Keep it short. &
    Outline how the metre is defined today.\\\hline
  \textbf{Describe} & Say what something is like or what happens. You do not need to say why.
    & Describe the pattern made by light through a narrow slit.\\\hline
  \textbf{Explain} & Say \emph{why} or \emph{how} it happens. Link cause and effect with
    \emph{because}, \emph{so}, \emph{therefore}. & Explain why the wave model cannot account
    for the threshold frequency.\\\hline
  \textbf{Compare} & Show how two things are the same \emph{and} how they are different. &
    Compare Newton's and Huygens' models of refraction.\\\hline
  \textbf{Analyse} & Break the topic into parts, show how the parts are linked, then say what
    this means. & Analyse how muons give evidence for special relativity.\\\hline
  \textbf{Assess} & Judge how good, how important or how large something is, with reasons. &
    Assess the validity of Galileo's lantern method.\\\hline
  \textbf{Evaluate} & Judge using criteria: strengths \emph{and} limitations. End with a clear
    judgement. & Evaluate the significance of polarisation for the model of light.\\\hline
  \textbf{Justify} & Support a statement or a conclusion with evidence and reasons. & Justify
    the claim that light is a transverse wave.\\\hline
  \textbf{Calculate} & Work out a number. Show the formula, the substitution with units, and
    the answer with a unit. & Calculate the stopping voltage for \qty{254}{nm} light on
    calcium.\\\hline
\end{longtable}}

You may also meet \emph{contrast} (give only the differences), \emph{discuss} (points for and
against) and \emph{recount} (retell the main events in order).

\textbf{How much to write.} 1--2 marks: one to three sentences. 3--4 marks: three or four
\emph{linked} points (cause $\to$ effect). 5--9 marks: a short plan, then one paragraph for
each idea. For \emph{assess}, \emph{evaluate} and \emph{justify}, finish with a judgement.
Use physics words (photon, work function, proper time), and draw a labelled diagram when it
helps.

\section*{Model answers from this module}

Each answer below would earn full marks. The marking criteria show how the marks are given,
in the NESA style: the top row is for full marks.

\subsection*{1. Outline (2 marks)}
\textbf{Question.} Outline how the metre is defined today.

\begin{modelanswer}
Since 1983, the speed of light in a vacuum has a fixed, exact value,
\qty{299792458}{m.s^{-1}}. One metre is the distance that light travels in a vacuum in
1/299\,792\,458 of a second. The second is measured with caesium atomic clocks.
\end{modelanswer}
\begin{markingcriteria}
  \criterion{States that $c$ has a fixed, exact value \emph{and} that the metre is the distance
    light travels in a vacuum in a fixed fraction of a second}{2}
  \criterion{States one of these points}{1}
\end{markingcriteria}

\subsection*{2. Explain (3 marks)}
\textbf{Question.} Explain why the wave model of light cannot account for the threshold
frequency in the photoelectric effect.

\begin{modelanswer}
In the wave model, the energy of light depends on its intensity (brightness), not on its
frequency. The energy arrives spread over the whole surface. So the wave model predicts that light
of \emph{any} frequency releases electrons if it is bright enough, or if we wait long enough.

But experiments show a threshold frequency $f_0$: below it, no electrons are released, however
bright the light. This contradicts the wave prediction, so the wave model fails. (In the photon
model, one photon gives all its energy, $hf$, to one electron, which escapes only if
$hf \ge \phi$.)
\end{modelanswer}
\begin{markingcriteria}
  \criterion{Links wave energy to intensity, gives the wave prediction and shows that the
    observed threshold contradicts it}{3}
  \criterion{Gives the wave prediction and the observation, with an incomplete link}{2}
  \criterion{States one relevant point}{1}
\end{markingcriteria}

\subsection*{3. Calculate (3 marks)}
\textbf{Question.} Ultraviolet light of wavelength \qty{254}{nm} falls on a calcium surface.
The work function of calcium is \qty{2.87}{eV}. Calculate the stopping voltage.

\begin{modelanswer}
Change the units first: $\lambda = \qty{254}{nm} = \qty{2.54e-7}{m}$ and
$\phi = 2.87 \times \qty{1.602e-19}{J} = \qty{4.598e-19}{J}$.\par
Then find the photon energy, use Einstein's equation, and use $K_\mathrm{max} = qV_0$:
\begin{align*}
  E &= \frac{hc}{\lambda} = \frac{\qty{6.626e-34}{J.s} \times \qty{3.00e8}{m.s^{-1}}}
       {\qty{2.54e-7}{m}} = \qty{7.826e-19}{J}\\
  K_\mathrm{max} &= hf - \phi = \qty{7.826e-19}{J} - \qty{4.598e-19}{J} = \qty{3.228e-19}{J}\\
  V_0 &= \frac{K_\mathrm{max}}{q} = \frac{\qty{3.228e-19}{J}}{\qty{1.602e-19}{C}}
       = \qty{2.015}{V}
\end{align*}
Answer: \final{$V_0 = \qty{2.02}{V}$} (3 significant figures). Check: in eV,
$K_\mathrm{max} = 4.885 - 2.87 = \qty{2.02}{eV}$, the same number as $V_0$ in volts \cmark.
\end{modelanswer}
\begin{markingcriteria}
  \criterion{Correct photon energy, with $\lambda$ in metres}{1}
  \criterion{Correct $K_\mathrm{max}$, with both energies in the same unit}{1}
  \criterion{Correct stopping voltage with its unit}{1}
\end{markingcriteria}

\subsection*{4. Analyse (5 marks)}
\textbf{Question.} Cosmic rays make muons \qty{10.0}{km} above sea level. The muons travel
straight down at $0.998c$ ($\qty{2.994e8}{m.s^{-1}}$). Their mean lifetime at rest is
\qty{2.20}{\micro\second}. Many of them are detected at sea level. Analyse how these
observations give evidence for special relativity. Include calculations.

\begin{modelanswer}
\textbf{The problem.} Without relativity, a muon travels
$vt_0 = \qty{2.994e8}{m.s^{-1}} \times \qty{2.20e-6}{s} = \qty{659}{m}$ in one lifetime. The
trip takes $\qty{1.00e4}{m} \div \qty{2.994e8}{m.s^{-1}} = \qty{33.4}{\micro\second}$, about
15 lifetimes. So almost no muons should reach sea level. But many do.

\textbf{Earth frame: time dilation.} The muon's own clock is present at its birth and its
decay, so \qty{2.20}{\micro\second} is the proper time $t_0$. Since
$\sqrt{1 - 0.998^2} = 0.0632$,
\[ t = \frac{t_0}{\sqrt{\left(1 - \dfrac{v^2}{c^2}\right)}} = \frac{\qty{2.20}{\micro\second}}{0.0632}
   = \qty{34.8}{\micro\second} \]
The trip (\qty{33.4}{\micro\second}) is shorter than one dilated lifetime, so many muons
survive.

\textbf{Muon frame: length contraction.} The atmosphere moves past the muon, so its thickness
is contracted:
\[ l = l_0\sqrt{\left(1 - \dfrac{v^2}{c^2}\right)} = \qty{1.00e4}{m} \times 0.0632 = \qty{632}{m} \]
This trip takes $\qty{632}{m} \div \qty{2.994e8}{m.s^{-1}} = \qty{2.11}{\micro\second}$, less
than the lifetime of \qty{2.20}{\micro\second}.

\textbf{What this means.} Both frames agree: the trip lasts about 0.96 of a mean lifetime.
Classical physics predicts almost no muons at sea level; relativity predicts many, as
observed. So muons are strong evidence for time dilation and length contraction, and so for
special relativity.
\end{modelanswer}
\begin{markingcriteria}
  \criterion{Shows why classical physics fails; uses correct calculations in \emph{both}
    frames; draws out what this means for special relativity}{5}
  \criterion{As above, with one frame or one calculation missing or wrong}{4}
  \criterion{Explains the observation with time dilation \emph{or} length contraction, with a
    correct calculation}{3}
  \criterion{Describes time dilation or length contraction and links it to muons}{2}
  \criterion{States one relevant point}{1}
\end{markingcriteria}

\subsection*{5. Evaluate (5 marks)}
\textbf{Question.} Evaluate the significance of polarisation in the development of a model of
light.

\begin{modelanswer}
Polarised light vibrates in one direction only. Only transverse waves can be polarised,
because only they have a sideways vibration for a filter to select.

\textbf{Strengths.} Polarisers (and Malus' law, $I = I_\mathrm{max}\cos^2\theta$) showed
that light can be polarised, so light must be a \emph{transverse} wave. Huygens' longitudinal
wave model had to change: Young and Fresnel proposed transverse light waves. Newton's particle
model had no natural explanation for polarisation. Later, Maxwell's theory predicted an EM
wave whose fields vibrate at \ang{90} to its direction of travel, which fits polarisation
exactly.

\textbf{Limitations.} On its own, polarisation does not show that light is electromagnetic,
and interference was also needed to show that light is a wave. Later, the photoelectric
effect showed that the wave model is not complete.

\textbf{Judgement.} Polarisation was highly significant: it decided what \emph{kind} of wave
light is and strongly supported Maxwell's EM wave model, although other evidence was needed
for the full model.
\end{modelanswer}
\begin{markingcriteria}
  \criterion{Clear judgement, supported by: light is transverse; the effect on Huygens' and
    Newton's models; the link to Maxwell; at least one limitation}{5}
  \criterion{Judgement supported by three of these points}{4}
  \criterion{Explains why polarisation shows light is transverse, and one effect on a model}{3}
  \criterion{Describes polarisation and states that light is transverse}{2}
  \criterion{States one relevant point}{1}
\end{markingcriteria}
